\documentclass[10pt, xetex]{article}
\linespread{0.94}
\usepackage[includeheadfoot,
            %showframe,% показывает поля --- нужно для отладки
            paperwidth=145mm, %ширина листа
            paperheight=210mm, %высота
            top=0.6cm, %верхние поля
            bottom=0.8 cm, %нижние          
            outer=10mm, %внешние
            inner=10mm, %внутренние
            marginparwidth=0.6cm, %ширина текста на полях
            marginparsep=1.2mm, %расстояние от тела текста до текста на полях 
            headheight=.5cm,
            footskip=0.6cm, %расстояние от нижней границы нижних полей до текста
            columnsep=.9cm,
            headsep=0.2cm, %расстояние от верхнего колонтитула до текста
            heightrounded
            ]{geometry}        
\setlength\parindent{8pt}
\usepackage[]{fontspec}
\usepackage[english, russian]{babel}
\setromanfont{STIXTwoText}
\setsansfont{STIXTwoText}
\setmonofont{Anonymous Pro}
\defaultfontfeatures{Ligatures={TeX}} 
\setmainfont[]{STIXTwoText}

\usepackage[math-style=TeX]{unicode-math}
\unimathsetup{math-style=TeX}
\setmathfont{texgyrepagella-math.otf}

\usepackage[framemethod=pstricks]{mdframed}
\usepackage{xkeyval}
\usepackage{ifthen}
\usepackage{ifpdf}
\usepackage{etoolbox}
\usepackage{multicol}

\usepackage{xcolor}

\usepackage{pstricks}
\usepackage{epsfig}
\usepackage{pst-circ}
\usepackage{pst-grad}
\usepackage{pst-plot}
\usepackage{pst-coil}
\usepackage{pst-slpe,pstricks-add,pst-blur, pst-func}
\newpsstyle{Blackball}{fillstyle=ccslope,slopebegin=white,slopeend=black,
    slopecenter={0.6 0.6}, linewidth=0.25pt, linecolor=black}


\usepackage[pstricks]{bclogo} \setlength{\logowidth}{10pt}
\renewcommand\bcStyleTitre[1]{\hskip1.6mm\emph{#1}}




\usepackage{titlesec}
\titlelabel{\thetitle.~}
\titleformat{name=\part}[hang]{\Huge}{}{}{\filcenter}{\filcenter}
\titlespacing*{\part}{0pt}{2pt plus 2pt minus 1pt}{10pt}%  
\titleformat{name=\section}[hang]{\Large}{}{2pt}{\filcenter\thesection.~}{\filcenter}
\titlespacing*{\section}{0pt}{8pt plus 2pt minus 1pt}{2pt}%   
\titleformat{name=\section,numberless}[hang]% уменьшение расстояния до section
   {\Large}%
   {}%
   {0em}%
   {\filcenter}%
\titlespacing{\section}{0pt}{8pt plus 2pt minus 1pt}{2pt}% 
\titleformat{name=\subsection,numberless}% уменьшение расстояния до subsection
   {\large}%
   {}
   {0em}%
   {\filcenter}%
\titlespacing*{\subsection}{0pc}{4pt plus 1pt minus 0.5pt}{0pt plus 0.5pt minus 0.5pt}%
 
%\usepackage{showframe}

\usepackage{graphicx}
\usepackage{fancybox,fancyhdr}
%\usepackage{eucal,bm, amsthm, amssymb, amsmath, array}
\usepackage[scale=1.18,boondox]{emf} 

\usepackage[hypcap=false]{caption}
\captionsetup[figure]{font={small,sl}, name={К примеру}}
\setlength{\abovecaptionskip}{3pt plus 1pt minus 1pt}
\captionsetup[table]{font={small,sl}, name={Таблица}}
\setlength{\abovecaptionskip}{3pt plus 1pt minus 1pt}

\usepackage{float}
\definecolor{darkred}{HTML}{8B0000}
\definecolor{linkcolor}{HTML}{120A8F}
%\definecolor{arrow_red}{HTML}{FF0000}

\setlength{\abovecaptionskip}{0.2cm}% от подписи до float 
\setlength{\textfloatsep}{5pt plus 1.0pt minus 2.0pt} %расстояние от float до текста;
\setlength{\floatsep}{2pt}%расстояние между двумя floats;
\setlength{\intextsep}{4pt}%расстояние между floats с [h] и текстом ;

\newcounter{zadacha}
\newcommand{\z}{\par\vspace{0.6\baselineskip plus 1pt minus 1pt}\noindent\addtocounter{zadacha}{1}\hypertarget{z\thezadacha}{}\hyperlink{an\thezadacha}{{Пример {\arabic{zadacha}.}\:\hskip0.15mm}}}
\newcounter{answer}
\newcommand{\an}{\par\vspace{0.25\baselineskip plus 1pt minus 1pt}\noindent\addtocounter{answer}{1}%
\hypertarget{an\theanswer}{}\hyperlink{z\theanswer}{\:\textbf{\arabic{answer}.}\:\hskip0.15mm}}


\renewcommand{\thefootnote}{\arabic{footnote}}
\usepackage[]{hyperref}
%\definecolor{linkcolor}{HTML}{1560BD} % цвет ссылок
\definecolor{violet}{HTML}{800080} % цвет гиперссылок
%\definecolor{arrow_blue}{HTML}{000080}%{000080}%{ADD8E6}}
\hypersetup{linkcolor=linkcolor,citecolor=darkred, urlcolor = violet, colorlinks=true}
\hypersetup{pdfhighlight=/P}
\usepackage[all]{hypcap}


\bibliographystyle{ugost2008l}


\begin{document}



\pagestyle{fancy}
\rhead{\small{Лето 2020}}
%\setlength{\headheight}{1cm}
\lhead{\uppercase{\small\textsl{Олимпиадный марафон СУНЦ, 8 кл.}}}
\cfoot{ --~\thepage~-- }
%\rfoot{\fontspec{Pecita}{Пётр А. Крюков}}
%\rfoot{\texttt{\textit{\small{Пётр А. Крюков}}}}
\renewcommand{\headrulewidth}{0.75pt} %линия в верхнем колонтитуле
\renewcommand{\footrulewidth}{0.0pt} %линия в нижнем колонтитуле
\setlength{\textfloatsep}{10pt plus 1.0pt minus 2.0pt}
\setlength{\floatsep}{6pt plus 1.0pt minus 1.0pt}
\setlength{\intextsep}{6pt plus 1.0pt minus 1.0pt}
\part*{Приближённые вычисления I}
\vskip-\baselineskip
\z Вычислите приближённо значения обыкновенных дробей: $\displaystyle{\frac{1}{3}}$, $\displaystyle{\frac{1}{7}}$, $\displaystyle{\frac{1}{9}}$, округлив до двух значащих цифр. На сколько процентов (с точностью до целых) в каждом случае отличается приближённое значение от точного?
\vskip0.25\baselineskip
Сделаем расчёт для $\displaystyle{\frac{1}{3}}$. Приближённое значение равно $0{,}33$, поэтому справедливо равенство 
\setlength{\belowdisplayskip}{0.4\baselineskip} 
\setlength{\abovedisplayskip}{0.4\baselineskip}
\begin{equation*}
\frac{\frac{1}{3}-0{,}33}{\frac{1}{3}} = 1 - 3\cdot 0{,}33 = 0{,}01 = 1\:\%.
\end{equation*}
Аналогичные вычисления для двух других дробей дают: $0{,}14$ и $2\:\%$; $0{,}11$ и $1\:\%$.

\z Число $\pi$ --- иррациональное. На сколько процентов приближённое значение числа $\pi$, равное $3$; $3{,}1$; $3{,}14$, отличается от приближённого значения, равного $3{,}14159$? Ответ дайте с точностью до одной значащей цифры. 
\vskip0.25\baselineskip
Расчёт для приближённого значения $3$. 
\setlength{\belowdisplayskip}{0.4\baselineskip} 
\setlength{\abovedisplayskip}{0.4\baselineskip}
\begin{equation*}
\frac{3{,}14159-3}{3{,}14159} \approx 0{,}045 = 4{,}5\:\%.
\end{equation*}
В двух других случаях получаются значения: $1\:\%$; $0{,}05\:\%$.

\z Радианная мера развёрнутого угла ($180^{\circ}$) равна $\pi$. Заполните пустые клетки таблицы \ref{angles_sins}. Используйте калькулятор. В этой задаче следует считать, что число $\pi$ приближённо равно $3{,}1416$. Тригонометрические функции следует вычислять от углов, заданных в градусах. Вычисленные значения следует округлить до десятитысячных. 
\vskip 0.2cm
\begin{table}[h]
\renewcommand{\arraystretch}{1.6} %% increase table row spacing
\renewcommand{\tabcolsep}{12 pt}
\small
\centering{
\begin{tabular}{|l|l|l|l|l|l|l|l|}
\hline
$\alpha$,~град & 0,5 & 1,0 & 2,0 & 5,0 & 10,0 & 15,0 & 30,0\\
\hline
$\alpha$,~рад &  &  &  &  &  &  &  \\
\hline
$\sin \alpha$ &  &  &  &  &  &  &  \\
\hline
$\tg \alpha$ &  &  &  &  &  &  &  \\
\hline
$\cos \alpha$ &  &  &  &  &  &  &  \\
\hline
\end{tabular}
}
\caption{}
\label{angles_sins}
\end{table}
\normalsize
\indent Угол в градусах пересчитывается в угол в радианах по формуле
\setlength{\belowdisplayskip}{0.4\baselineskip} 
\setlength{\abovedisplayskip}{0.4\baselineskip}
\begin{equation*}
\alpha\:(\text{рад}) = \frac{\alpha\:(\text{град})}{180}\cdot \pi.
\end{equation*}

\z Определите значения угла $\alpha$ в градусах (округлив до десятых) и радианах (округлив до сотых), если известен $\sin\alpha$, таблица \ref{arcsin}. Если ваш калькулятор снабжён функцией $\arcsin$ (обычно обозначается <<asin>>), возвращающей значение угла для заданного синуса, можно её использовать. Если нет, рекомендуем решать подбором, используя функцию sin.
\vskip 0.2cm
\begin{table}[h]
\renewcommand{\arraystretch}{1.6} %% increase table row spacing
\renewcommand{\tabcolsep}{12 pt}
\small
\centering{
\begin{tabular}{|l|l|l|l|l|l|}
\hline
$\sin \alpha$& 0,2 & 0,3 & 0,5 & 0,75 & 0,9\\
\hline
$\alpha$,~град  &  &  &  &  &      \\
\hline
$\alpha$,~рад&  &  &  &  &    \\
\hline
\end{tabular}
}
\caption{}
\label{arcsin}
\end{table}


\vskip-0.5\baselineskip
\begin{bclogo}[logo=\bccrayon, margeG=-0.4, margeD=0, noborder = true, couleurBarre = darkred]{Погрешность (или точность) приближения}
Пусть $f$ --- точное значение некоторой величины, а $f_{a}$ --- приближённое, обозначим: $\Delta f = |f-f_{a}|$ --- \textsl{отклонение} приближённого значения от точного. Будем называть  \textsl{погрешностью приближения} отношение: $\delta f =\displaystyle{\frac{\Delta f}{f}}$. Договоримся выражать погрешность приближения в процентах, округляя до одной или двух значащих цифр. Словосочетание \textsl{точность приближения} имеет тот же смысл, что и \textsl{погрешность приближения}.
\end{bclogo}
\vskip-0.5\baselineskip

\z При каких $x$ в выражении $\displaystyle{x+\frac{1}{x}}$ можно пренебречь первым слагаемым, а при каких вторым, так чтобы точность приближения была не хуже $4\:\%$ в каждом случае?
\vskip0.25\baselineskip
Обозначим: $f=\displaystyle{x+\frac{1}{x}}$. Если первым слагаемым можно пренебречь, то:  $x \ll \displaystyle{\frac{1}{x}}$,  $f_{a} = \displaystyle{\frac{1}{x}}$, для определённости считаем $x>0$. Тогда:
\begin{equation*}
\delta f = \frac{x}{x+\cfrac{1}{x}} = \frac{x^{2}}{x^{2}+1}\approx x^{2},
\label{example_1_f}
\end{equation*}
откуда следует, что при $x$: $x^{2}< 0{,}04$, $f \approx \displaystyle{\frac{1}{x}}$ c точностью не хуже $4\:\%$. Легко видеть, что: $x < 0{,}2$.\\ 
\indent Во втором случае, когда можно пренебречь слагаемым $\displaystyle{\frac{1}{x}}$ по сравнению с $x$, делая аналогично, получаем: $\delta f \approx \displaystyle{\frac{1}{x^{2}}}$, $x> 5$.
\setcounter{figure}{\thezadacha}
\z Пусть $P_{4}$ --- периметр квадрата, описанного около окружности радиусом $R$, а $p_4$ --- периметр квадрата, вписанного в эту окружность (см. рис. \ref{multi_angl}). Полусумма периметров $L_{4} =\displaystyle{\frac{P_{4}+p_4}{2}}$ даёт оценку длины окружности. Аналогичные оценки можно сделать, используя вместо квадрата любые другие правильные $n$-угольники. Точность приближения при этом можно оценить (!) по формуле $\delta L_{n} \approx \displaystyle{\frac{P_{n}-p_{n}}{P_{n}+p_n}}$. Определите точность $\delta L_{n}$ для $n=4$ и $n = 6$.

\begin{figure}[h]
\centering
{
\includegraphics[]{multi_angl.eps}
}
\caption{}
       \label{multi_angl} 
\end{figure}

Решим для шестиугольников. Длина стороны вписанного шестиугольника равна радиусу $R$, длина стороны описанного шестиугольника находится по теореме Пифагора: $\displaystyle{\frac{2R}{\sqrt{3}}}$. По формуле из условия:
\begin{equation*}
\delta L_{6} = \frac{\frac{2}{\sqrt{3}}-1}{\frac{2}{\sqrt{3}}-1} = \frac{2-\sqrt{3}}{2+\sqrt{3}}\approx 7\:\%.
\end{equation*}
Отметим, что уже для $n=12$ получится $\delta L_{12}= 0{,}14\:\%$.

\z Округление физических констант
\begin{table}[h]
\renewcommand{\arraystretch}{1.6} %% increase table row spacing
\renewcommand{\tabcolsep}{10 pt}
\centering{
\begin{tabular}{|l|c|c|c|}
\hline
Постоянная & Таблица & Округление & Точность  \\
\hline
G,\:$\cdot10^{-11}$\:$\hbox{Н}\cdot\hbox{м}^{2}/\hbox{кг}^{2}$ &6,67&7  &5\%     \\
\hline
$g$,\:$\hbox{Н}/\hbox{кг}$  &9,8& 10&2\%      \\
\hline
$R$,\:$\hbox{Дж}/(\hbox{моль}\cdot\hbox{К}) $  &8,31& 8&3,7\%      \\
\hline
$RT_{(t = 20{^{\circ}}\,{\textup{C})}}$,\:$\hbox{Дж}/\hbox{моль}$  & 2434,8&$2500$ & 2,7\%   \\
\hline
$c_{\hbox{\footnotesize{вод.}}}$,\:$\hbox{Дж}/(\hbox{кг}\,{^{\circ}}\,{\hbox{C}})$  &4180 &$4000$ & 4,3\%   \\
\hline
\end{tabular}
}
\caption{}
\end{table}
\z В пачке бумаги $500 \pm 10$ листов. Толщину пачки измеряют, оценивают погрешность, получают:
$D = 45 \pm 1~\hbox{мм}$. Чему равна наиболее вероятная толщина одного листа $d_{0}$? Оцените точность определения толщины листа $\delta d$.
\vskip0.25\baselineskip
Cлова в условии <<\textsl{наиболее вероятная толщина}>> --- дань терминологии. Найдём $d_{\text{min}}$ и $d_{\text{max}}$. Имеем: $d_{\text{min}} = \displaystyle{\frac{44~\hbox{мм}}{510} \approx 8{,}6\cdot10^{-2}~\hbox{мм}}$, $d_{\text{max}} \approx 9{,}4\cdot10^{-2}~\hbox{мм}.$ Далее ответ: $d_{0}=\displaystyle{\frac{d_{\text{min}}+d_{\text{max}}}{2} \approx 9{,}0\cdot10^{-2}~\hbox{мм}}$, $\delta d =\displaystyle{\frac{d_{\text{max}}-d_{\text{min}}}{d_{\text{max}}+d_{\text{min}}} = 4\:\%}$.
\vskip-0.5\baselineskip
\begin{bclogo}[logo=\bccrayon, margeG=-0.4, margeD=0, noborder = true, couleurBarre = darkred]{Погрешность суммы и разности}
Пусть приближённое значение величины $x$ равно $x_{a}$ и оценка отклонения равна  $\Delta x$, а приближённое значение величины $y$ равно $y_{a}$ с отклонением $\Delta y$, тогда приближённое значение величины $f = x\pm y$ равно $f_{a} = x_{a} \pm y_{a}$, а  отклонение приближённого значения от точного имеет оценку $\Delta f = \Delta x + \Delta y$.
\end{bclogo}
\vskip-0.5\baselineskip

\z Пусть $S = x+y$, $D = x-y$. Известно, что: $x = 10 \pm 1$, $y = 8 \pm 1$. Определите величину $\delta S$ и $\delta D$.
\vskip0.25\baselineskip
Отметим, что точности заданных величин равны: $\delta x = 10\%$, $\delta y = 12{,}5\%$.\\
По определению отклонения суммы и разности равны соответственно: $\Delta S =$ $= \Delta D = 2$. \\
Точность определения суммы равна $\displaystyle{\delta S =\frac{\Delta S}{S} = \frac{2}{18}\approx 11\:\%}$.\\
Точность определения разности равна $\displaystyle{\delta D =\frac{\Delta D}{D} = \frac{2}{2}=100\:\%}$.
Получается, что по данным в условии $x$ и $y$ определить $x-y$ нельзя!

\vskip-0.05\baselineskip
\begin{bclogo}[logo=\bccrayon, margeG=-0.4, margeD=0, noborder = true, couleurBarre = darkred]{Оценка точности произведения и частного}
Пусть $x_{a}$ --- приближённое значение величины $x$, найденное с точностью $\delta x$, при этом $\delta x << 1$, например $\delta x < 10\%$. Пусть $y_{a}$ --- приближённое значение величины $y$, найденное с точностью $\delta y$, $\delta y << 1$. Тогда приближённое значение величины $M = x\cdot y$ равно $M_{a} = x_{a}y_{a}$, а его точность имеет оценку $\delta M = \delta x + \delta y$. Приближённое значение величины $D = \displaystyle{\frac{x}{y}}$ равно $D_{a} = \displaystyle{\frac{x_{a}}{y_{a}}}$ с точностью $\delta D = \delta x + \delta y$. 
\end{bclogo}
\vskip-0.75\baselineskip

\z Приближённое значение радиуса $R$ окружности равно $R_{a}$ и найдено с точностью $5\:\%$. Приближённое значение площади соответствующего круга определяют по формуле: $S_{a} \approx 3R_{a}^{2}$. Найдите точность приближения. Какой наилучшей точности определения $S$ можно достичь при заданной точности $\delta R$?
\vskip0.25\baselineskip
Поскольку $\pi \approx 3$ c точностью $\delta \pi = 5\%$ искомая точность равна 
\setlength{\belowdisplayskip}{0.25\baselineskip} 
\setlength{\abovedisplayskip}{0.25\baselineskip}
\begin{equation*}
\delta S = \delta \pi + 2\cdot \delta R = 15\%.
\end{equation*}
\begin{equation*}
\delta S_{\hbox{\footnotesize{нилучш.}}} = 2\cdot \delta R = 10\%.
\end{equation*}

\z В одной  задаче по графику определялись значения моментов времени: $T_{1}=3{,}01\pm 0{,}03~\hbox{с}$, $T_{2}=3{,}89\pm 0{,}03~\hbox{с}$, $T_{3}= 4{,}61\pm 0{,}03~\hbox{с}$. Далее вычисляли величины: $t_{1} = T_{2}- T_{1}$ и $t_{2} = T_{3}- T_{2}$. Далее определяли значение коэффициента $k = \displaystyle{\frac{t_{2}}{t_{1}}}$ и наконец по формуле 
\setlength{\belowdisplayskip}{0.25\baselineskip} 
\setlength{\abovedisplayskip}{0.25\baselineskip}
\begin{equation*}
H =\frac{gt_{1}^{2}}{8 k^{2}},
\end{equation*}
где $g$ --- ускорение свободного падения, определяли искомую величину $H$. Чему равна точность определения $H$ по данным задачи?
\vskip0.25\baselineskip
Найдём сначала точность приближённых значений $\delta t_{1}$ и $\delta t_{2}$ 
\setlength{\belowdisplayskip}{0.25\baselineskip} 
\setlength{\abovedisplayskip}{0.25\baselineskip}
\begin{equation*}
\delta t_{1} = \frac{0{,}06}{0{,}88} \approx 7\,\%, \quad \delta t_{2} = \frac{0{,}06}{0{,}72} \approx 8\,\%.
\end{equation*}
Точность вычисления коэффициента $k$ равна $\delta k \approx 15\,\% $. Таким образом в лучшем случае искомая точность равна 
\setlength{\belowdisplayskip}{0.25\baselineskip} 
\setlength{\abovedisplayskip}{0.25\baselineskip}
\begin{equation*}
\delta H = 2\cdot\delta t_{1}+2\cdot\delta k = 44\,\%.
\end{equation*}
Это называется \textsl{накопление ошибок}.

\vskip-0.5\baselineskip
\begin{bclogo}[logo=\bccrayon, margeG=-0.4, margeD=0, noborder = true, couleurBarre = darkred]{Приближённые формулы}%
Для малых углов $\alpha$ (если $\alpha$ выражен в радианах, то $|\alpha|\ll 1$) справедливы соотношения:
\setlength{\belowdisplayskip}{0.5\baselineskip} 
\setlength{\abovedisplayskip}{0.4\baselineskip}
\begin{equation}
\sin \alpha \approx \alpha, \quad  \tg \alpha \approx \alpha, \quad \cos \alpha \approx 1.
\label{sin_tan_approx}
\end{equation}
Для $|x| \ll 1$ верны приближённые равенства:
\begin{equation}
\frac{1}{1+x} \approx 1-x, \quad  (1+x)^{2} \approx 1+2x, \quad \sqrt{1+x} \approx 1+\frac{x}{2}.
\label{power_approx}
\end{equation} 
\end{bclogo}
\vskip-0.5\baselineskip
\z Для приближённых формул из первого столбца таблицы \ref{polynom_est} впишите в свободные клетки таблицы точность приближения для разных $x$.
\vskip 0.05cm
\begin{table}[h]
\renewcommand{\arraystretch}{1.8} %% increase table row spacing
\renewcommand{\tabcolsep}{12 pt}

\small
\centering{
\begin{tabular}{|l|l|l|l|l|}
\hline
  $|x|$ & 0,01 & 0,10 & 0,20 & 0,40  \\
\hline
\raisebox{0.2\height}{$\frac{1}{1+x}\approx 1-x$}&  &   &  &  \\
\hline
\raisebox{0.2\height}{$\sqrt{1+x}\approx 1+\frac{x}{2}$} &   &   &   &   \\
\hline
\raisebox{0.2\height}{$(1+x)^{2}\approx 1+2x$} & &  &  &  \\
\hline
\end{tabular}
}
\caption{}
\label{polynom_est}
\end{table}
\vskip-0.5\baselineskip
\z Точность приближения для тригонометрических функций
\begin{table}[h]
\renewcommand{\arraystretch}{1.5} %% increase table row spacing
\renewcommand{\tabcolsep}{8 pt}
\small
\centering{
\begin{tabular}{|l|c|c|c|c|c|}
\hline
$\alpha$,~град &  1,0 &  5,0 & 10,0 & 15,0 & 30,0\\
\hline
$\alpha$,~рад & 0,0175  &0,0873  &0,1745  &0,2618  &0,5236  \\
\hline
$\sin \alpha$ &0,0175  &0,0872  &0,1736  &0,2588  &0,5000  \\
\hline
$\delta (\sin)$ & $< 0{,}1\:\%$  &$0{,}1\:\%$  &$0{,}5\:\%$  &$1{,}1\:\%$  &$4{,}5\:\%$  \\
\hline
$\tg \alpha$ &0,0175  &0,0875  &0,1763  &0,2679  &0,5774  \\
\hline
$\delta (\tg)$ & $< 0{,}1\:\%$  &$0{,}2\:\%$  &$1\:\%$  &$2{,}3\:\%$  &$9{,}3\:\%$  \\
\hline
\end{tabular}
}
\caption{}
\label{sin_tan_est}
\end{table}
\z Если $(1-x)$ --- частное от деления $1$ на $(1+x)$, то чему может быть равен остаток? Получите выражение в виде многочлена от $x$. 
\vskip0.25\baselineskip
 Ниже показано деление единицы на $1+x$, которое и даёт ответ к задаче.  

\[
\arraycolsep=0.05em
\renewcommand{\arraystretch}{1.1}
\begin{array}{rrrrr@{\,}r|l}
1&&&&&& {\,}1+x\\
\cline{7-7}
1&+x&&&&&   {\,}{\vcenter{\hbox{$1-x+x^{2}$}}}\\
\cline{1-2}
&-x&  \\
&-x& -x^{2} \\
\cline{2-3}
& & {\vcenter{\hbox{$x^{2}$}}}\\
& & x^{2}&  +x^{3}\\
\cline{3-4}
& & &  {\vcenter{\hbox{$-x^{3}$}}}
\end{array}
\]
Таким образом, если частное равно $1-x$, то остаток равен $x^{2}$. Нелинейный член в частном показан, чтобы стало понятно, что процесс можно продолжать до бесконечности. 

\vskip-0.25\baselineskip
\begin{bclogo}[logo=\bccrayon, margeG=-0.4, margeD=0, noborder = true, couleurBarre = darkred]{Оценка с точностью до $x^{n}$}
Будем говорить, что функция $f(x)$ при малых $x$ \textit{имеет оценку с точностью до} $x^{n}$, если она представляется (при малых $x$) в виде: 
\begin{equation*}
f(x) \approx f(0)+a_{1}x+a_{2}x^{2}+\ldots+a_{n}x^{n},
\label{small_x_approx}
\end{equation*}
где $f(0)$ --- значение функции $f(x)$ при $x=0$; $a_{1}, a_{2},\ldots,a_{n}$ --- действительные числа. Оценка с точностью до $x$ называется \textit{линейным приближением}.
\end{bclogo}
\vskip-0.5\baselineskip
\z Из предыдущего примера видно, что с точностью до $x^{2}$ при малых $x$ справедливо приближённое равенство
\[
\frac{1}{1+x} \approx 1-x+x^{2}.
\]
Также становится ясно, как получить более точные приближения. 

\z Оценивая в линейном по малому параметру $x$ приближении алгебраические выражения, следует на каждом шаге пренебрегать слагаемыми, пропорциональными старшим степеням $x$. Например:

\[
(1+x)(2+x)(3+x)= 6 + 11x+6x^{2}+x^{3}\approx 6 + 11x;
\]
\[
\frac{1-x}{2+3x} = \frac{1-x}{2}\cdot \frac{1}{1+\frac{3x}{2}}\approx \frac{1}{2}(1-x)\left(1-\frac{3x}{2}\right)\approx \frac{1}{2}\left(1-\frac{5x}{2}\right);
\]
\[
\frac{1-x}{1+x+x^{2}} \approx \frac{1-x}{1+x}\approx (1-x)^{2}\approx 1-2x.
\]
\z Используя условия $\Delta r \ll r$, $\Delta r \ll R$ преобразуйте выражение:
 \[
\frac{R(r + \Delta r)}{R+r + \Delta r} = \frac{Rr}{R+r}\cdot \frac{1+\frac{\Delta r}{r}}{1+\frac{\Delta r}{R+r}}\approx \frac{Rr}{R+r}\left(1+\frac{\Delta r}{r}-\frac{\Delta r}{R+r} \right)=\frac{Rr}{R+r}\left(1+\delta r\frac{R}{R+r} \right).
\]
Введено обозначение малого параметра: $\delta r = \displaystyle{\frac{\Delta r}{r}}$.

\z Оцените в линейном приближении при малых $x$ выражение:
\[
\frac{\sqrt{(x+2)(3x+2)}}{3+5x}.
\]
\vskip0.25\baselineskip
Преобразуем сначала числитель, используя формулу для корня из (\ref{power_approx}). 
\[
\sqrt{(x+2)(3x+2)} \approx \sqrt{4+8x} = 2\sqrt{1+2x}\approx 2\left(1+\frac{2x}{2}\right) = 2(1+x).
\]
Исходное выражение теперь принимает вид:
\[
\frac{\sqrt{(x+2)(3x+2)}}{3+5x}\approx \frac{2(1+x)}{3+5x} = \frac{2}{3}\cdot \frac{1+x}{1+\frac{5x}{3}}.
\]
<<Переносим>> знаменатель в числитель и получаем ответ:
\[
\frac{2}{3}\cdot \frac{1+x}{1+\frac{5x}{3}} \approx \frac{2}{3}(1+x)\left(1-\frac{5x}{3}\right) \approx \frac{2}{3}\left(1-\frac{2x}{3}\right).
\]

\z Приближённые формулы $\sin x \approx x$ и $\tg x \approx x$ справедливы не только в линейном, но и в квадратичном (с точностью до $x^{2}$) приближении, а формула $\cos x \approx 1$ верна только в линейном приближении. Получите приближённую формулу для $\cos x$ в квадратичном приближении. 
\vskip0.25\baselineskip
Рассмотрим прямоугольный треугольник с гипотенузой, равной $1$, и острым углом, равным $x$. Один из его катетов равен $\cos x$, а другой равен $\sin x \approx x$ с точностью до слагаемых, пропорциональных $x^{2}$. По теореме Пифагора:
\[
\cos^{2}x \approx 1-x^{2},
\]
Откуда, используя приближённую формулу для корня, получаем ответ:
\[
\cos x \approx 1-\frac{x^{2}}{2}.
\]



\end{document}